% Options for packages loaded elsewhere
\PassOptionsToPackage{unicode}{hyperref}
\PassOptionsToPackage{hyphens}{url}
\documentclass[
  ignorenonframetext,
]{beamer}
\newif\ifbibliography
\usepackage{pgfpages}
\setbeamertemplate{caption}[numbered]
\setbeamertemplate{caption label separator}{: }
\setbeamercolor{caption name}{fg=normal text.fg}
\beamertemplatenavigationsymbolsempty
% remove section numbering
\setbeamertemplate{part page}{
  \centering
  \begin{beamercolorbox}[sep=16pt,center]{part title}
    \usebeamerfont{part title}\insertpart\par
  \end{beamercolorbox}
}
\setbeamertemplate{section page}{
  \centering
  \begin{beamercolorbox}[sep=12pt,center]{section title}
    \usebeamerfont{section title}\insertsection\par
  \end{beamercolorbox}
}
\setbeamertemplate{subsection page}{
  \centering
  \begin{beamercolorbox}[sep=8pt,center]{subsection title}
    \usebeamerfont{subsection title}\insertsubsection\par
  \end{beamercolorbox}
}
% Prevent slide breaks in the middle of a paragraph
\widowpenalties 1 10000
\raggedbottom
\AtBeginPart{
  \frame{\partpage}
}
\AtBeginSection{
  \ifbibliography
  \else
    \frame{\sectionpage}
  \fi
}
\AtBeginSubsection{
  \frame{\subsectionpage}
}
\usepackage{iftex}
\ifPDFTeX
  \usepackage[T1]{fontenc}
  \usepackage[utf8]{inputenc}
  \usepackage{textcomp} % provide euro and other symbols
\else % if luatex or xetex
  \usepackage{unicode-math} % this also loads fontspec
  \defaultfontfeatures{Scale=MatchLowercase}
  \defaultfontfeatures[\rmfamily]{Ligatures=TeX,Scale=1}
\fi
\usepackage{lmodern}
\ifPDFTeX\else
  % xetex/luatex font selection
\fi
% Use upquote if available, for straight quotes in verbatim environments
\IfFileExists{upquote.sty}{\usepackage{upquote}}{}
\IfFileExists{microtype.sty}{% use microtype if available
  \usepackage[]{microtype}
  \UseMicrotypeSet[protrusion]{basicmath} % disable protrusion for tt fonts
}{}
\makeatletter
\@ifundefined{KOMAClassName}{% if non-KOMA class
  \IfFileExists{parskip.sty}{%
    \usepackage{parskip}
  }{% else
    \setlength{\parindent}{0pt}
    \setlength{\parskip}{6pt plus 2pt minus 1pt}}
}{% if KOMA class
  \KOMAoptions{parskip=half}}
\makeatother
\usepackage{longtable,booktabs,array}
\newcounter{none} % for unnumbered tables
\usepackage{calc} % for calculating minipage widths
\usepackage{caption}
% Make caption package work with longtable
\makeatletter
\def\fnum@table{\tablename~\thetable}
\makeatother
\setlength{\emergencystretch}{3em} % prevent overfull lines
\providecommand{\tightlist}{%
  \setlength{\itemsep}{0pt}\setlength{\parskip}{0pt}}
% Slide layout defaults for warning-clean scientific decks.
\usepackage{etoolbox}
\IfFileExists{xurl.sty}{\usepackage{xurl}}{}
\IfFileExists{seqsplit.sty}{\usepackage{seqsplit}}{\newcommand{\seqsplit}[1]{#1}}
\protected\def\breakseq#1{\seqsplit{#1}}
\protected\def\breaktt#1{\begingroup\ttfamily\seqsplit{#1}\endgroup}
\setlength{\emergencystretch}{6em}
\tolerance=5000
\hbadness=10000
\hfuzz=1pt
\setlength{\tabcolsep}{2pt}
\AtBeginEnvironment{longtable}{\tiny\renewcommand{\arraystretch}{0.86}\setlength{\tabcolsep}{1pt}}
\AtBeginEnvironment{tabular}{\tiny\renewcommand{\arraystretch}{0.86}\setlength{\tabcolsep}{1pt}}

% Natbib and cross-reference fallbacks — slides don't load natbib
% or cleveref, but combined-PDF manuscript prose may emit these
% commands. The fallback renders citations as a bracketed key list
% and cross-references as detokenized labels so slides don't fail on
% undefined control sequences or raw underscores. \providecommand is
% a no-op if packages load later.
\providecommand{\citep}[1]{[#1]}
\providecommand{\citet}[1]{#1}
\providecommand{\citealp}[1]{#1}
\providecommand{\citeauthor}[1]{#1}
\providecommand{\citeyear}[1]{#1}
\providecommand{\cref}[1]{\texttt{\detokenize{#1}}}
\providecommand{\Cref}[1]{\texttt{\detokenize{#1}}}

\newtheorem{proposition}{Proposition}
\newtheorem{hypothesis}{Hypothesis}
\usepackage{bookmark}
\IfFileExists{xurl.sty}{\usepackage{xurl}}{} % add URL line breaks if available
\urlstyle{same}
% fallback for those not using the hyperref driver hyperxmp:
\makeatletter
\@ifundefined{xmpquote}{\newcommand{\xmpquote}[1]{#1}}{}
\makeatother
\hypersetup{
  hidelinks,
  pdfcreator={LaTeX via pandoc}}

\author{\texorpdfstring{}{}}
\date{}

\begin{document}

\section{Introduction}\label{introduction}

\begin{frame}[allowframebreaks]{Introduction}
Disclosure control---the process of sanitizing classified or sensitive
information before public release---requires multiple layers of
validation to ensure that no source identities, operational details,
time-place selectors, or controlled dissemination markers leak into the
public record. This exemplar implements a reproducible pipeline that
combines text-level disclosure control with visual redaction proofing,
steganographic provenance embedding, and hardware-backed TPM sealing.

The pipeline operates on invented fixture data: synthetic segments with
classification levels ranging from UNCLASSIFIED through
TOP\_SECRET//SCI, synthetic redaction decisions with bounded reasons,
and synthetic reviewer records. No real source material is used. The
exemplar demonstrates the full release-review workflow:
classification-ceiling enforcement, redaction-span validation,
source-control coverage, mosaic-risk scoring, residual-marker detection,
review-gate evaluation, source-safe ledger generation, and TPM-sealed
sidecar production.
\end{frame}

\begin{frame}[allowframebreaks]{Scope and Safety Boundaries}
\protect\phantomsection\label{scope-and-safety-boundaries}
This exemplar is limited to lawful redaction, declassification support,
public-records release review, taxonomy normalization, source-safe
ledgers, reviewer approval gates, sanitized packet export, and
source-protection auditing. It does not provide targeting, collection,
evasion, or surveillance operational guidance. All committed examples
are invented fixtures.
\end{frame}

\begin{frame}[fragile,allowframebreaks]{Contributions}
\protect\phantomsection\label{contributions}
\begin{enumerate}
\tightlist
\item
  A classification taxonomy with SCI alias support and configurable
  public ceilings.
\item
  A release-audit engine that validates redaction spans, checks orphan
  decisions, enforces source-control coverage, and scores mosaic risk.
\item
  A source-safe redaction ledger that records SHA-256 hashes of redacted
  spans without exposing source text.
\item
  A visual proof matrix that enumerates four redaction styles across
  four PDF backgrounds, producing sixteen variant PDFs with identical
  source-safe decisions.
\item
  A steganography layer that post-processes each base PDF with nine
  security methods, producing provenance-enhanced companion PDFs with
  hash manifests.
\item
  Optional Kmyth TPM sealing that wraps each hash manifest and
  steganography PDF in \texttt{.ski} sidecars sealed against the
  TPM2-TSS storage hierarchy.
\item
  A comprehensive release packet that combines sanitized text, audit
  findings, ledger, hashes, review gate status, and paragraph audit
  tables.
\end{enumerate}
\end{frame}

\section{Architecture}\label{architecture}

\begin{frame}[allowframebreaks]{Architecture}
The pipeline architecture consists of two layers. Layer one performs
text-level disclosure control: each segment is audited against a public
classification ceiling, redaction spans are validated for non-overlap,
orphan decisions are flagged, and source-control coverage is enforced.
Layer two applies visual redaction treatments, steganographic provenance
overlays, and optional Kmyth TPM sidecar sealing.
\end{frame}

\begin{frame}[fragile,allowframebreaks]{Layer One: Disclosure Control}
\protect\phantomsection\label{layer-one-disclosure-control}
The disclosure-control engine operates on \breaktt{RedactionSegment}
objects, each carrying an identifier, classification level, text
content, and optional source controls. Redaction decisions are
\breaktt{RedactionDecision} records with bounded reasons:
\texttt{source\_identity}, \breaktt{operational\_detail},
\breaktt{time\_place\_selector}, \texttt{legal\_privilege}, and
\texttt{privacy}. The audit engine validates that:

\begin{itemize}
\tightlist
\item
  Redaction spans are non-overlapping and within text bounds.
\item
  Each segment above the public ceiling has at least one redaction
  decision.
\item
  Each source control has a corresponding \texttt{source\_identity}
  redaction.
\item
  No orphan decisions reference missing segments.
\item
  Residual markers (email, IP, coordinates, NOFORN, HUMINT, SIGINT) are
  absent from sanitized text.
\item
  The mosaic risk score---residual markers normalized by segment and
  pattern count---does not exceed the policy threshold.
\end{itemize}
\end{frame}

\begin{frame}[allowframebreaks]{Layer Two: Visual Proof and
Steganography}
\protect\phantomsection\label{layer-two-visual-proof-and-steganography}
Visual proofing is parameterized separately from the release-audit text
path. Four redaction styles---blackout (solid black fill), whiteout
(white fill with gray text), grayout (mid-gray fill), and blur
(offset-rendered token)---are rendered across four PDF
backgrounds---white, gray, black, and blur (subdued text with blur
effect). The 4×4 matrix yields sixteen variant PDFs, each receiving the
same source-safe redaction decisions.

The steganography layer post-processes each base PDF with nine security
methods:

{\def\LTcaptype{none} % do not increment counter
\begin{longtable}[]{@{}
  >{\raggedright\arraybackslash}p{(\linewidth - 2\tabcolsep) * \real{0.4706}}
  >{\raggedright\arraybackslash}p{(\linewidth - 2\tabcolsep) * \real{0.5294}}@{}}
\toprule\noalign{}
\begin{minipage}[b]{\linewidth}\raggedright
Method
\end{minipage} & \begin{minipage}[b]{\linewidth}\raggedright
Purpose
\end{minipage} \\
\midrule\noalign{}
\endhead
SHA-256/SHA-512 hash manifest & Cryptographic integrity verification \\
Diagonal watermark overlay & Visible provenance stamp \\
Footer provenance overlay & Page-level release metadata \\
First-page invisible text & Hidden identification marker \\
QR payload barcode & Machine-readable provenance \\
Code128 page barcode & Per-page tracking barcode \\
PDF Info metadata & Document-level metadata injection \\
XMP metadata & Standards-compliant metadata embedding \\
Embedded stego manifest attachment & Self-contained provenance
archive \\
\bottomrule\noalign{}
\end{longtable}
}
\end{frame}

\begin{frame}[fragile,allowframebreaks]{Layer Three: Kmyth TPM Sealing}
\protect\phantomsection\label{layer-three-kmyth-tpm-sealing}
Kmyth TPM sealing wraps each hash manifest and steganography PDF in a
\texttt{.ski} sidecar sealed against the TPM2-TSS storage hierarchy. The
sealed objects are bound to PCR selections and policy or-values,
ensuring that unsealing requires the same platform configuration.

On macOS, which lacks a hardware TPM, a software TPM emulator (swtpm)
and an mssim-to-swtpm protocol proxy bridge the TPM2-TSS mssim TCTI to
swtpm's native socket protocol. The proxy:

\begin{itemize}
\tightlist
\item
  \textbf{Control channel}: intercepts MS simulator platform commands
  (POWER\_ON=1, NV\_ON=11, TPM\_SESSION\_END=20) and returns success,
  since swtpm's \breaktt{-\/-flags\ startup-clear} handles TPM
  initialization.
\item
  \textbf{Data channel}: strips the 9-byte mssim wire header (4B
  cmd\_type + 1B locality + 4B tpm\_size), forwards raw TPM commands to
  swtpm, and wraps responses back in mssim format.
\end{itemize}

The kmyth-seal binary was patched to call
\breaktt{Tss2\_Sys\_FlushContext} for the storage key handle before
freeing TPM2 resources. Without this patch, consecutive kmyth-seal
invocations exhaust the swtpm transient object slots (typically three),
causing ``out of memory for object contexts'' errors on the second seal.
\end{frame}

\end{document}
